\chapter{Approximation Theory for Reduced Order Models} \label{chapter: results}




This chapter provides the main results of this thesis based on the data-driven ROMs 
introduced in Chapter (\ref{chapter: problem}). As such we repeat, that in order to
execute the reduction for the norm $\norm{\cdot}$ defined in that chapter,
we require a mass matrix, that is independent of $\mu \in \Theta$, since otherwise the orthonormality 
with regard to it cannot be computed under the generalized POD Algorithm (\ref{algo: pod}). However, we note, 
that the paper \cite{brivio2023} we base most of this chapter on, and severall other works dealing with a similar setup, instead
neglect the representation in the underlying Hilbert space and use a simple euklidian norm for the
coefficient vectors. For our sitution though, we let
\begin{equation} \label{eq: representation uhat}
    \hat{u}_h(x, t; \mu, \nu) = \sum_{j=1}^{N_h} [\uhhat(t)]_j \varphi_j(x),
\end{equation}  
for all $x \in \Omega$, $t \in [0, T]$ and parameters $(\mu, \nu) \in \Theta \times \Theta'$, which simultaneously 
specifies $[\G]_{ij} = \langle \varphi_i, \varphi_j \rangle = [M(\mu)]_{ij}$ as the mass matrix. The norm for
this chapter will be, if not stated otherwise, the norm from Equation (\ref{eq: definition mass matrix norm}), so
here
\begin{equation*}
    \norm{u} := \sqrt{u^T \G u},
\end{equation*} 
for all $u \in \R^{N_h}$.
Furthermore
the solution vector $\uhhat(t) = \uhthat \in \R^{N_h}$ is set to be constant between the time instants 
in our model, i.e. for all
$t \in [(k-1) \Delta t, k \Delta t)$ for some $k = 1, \dots, N_t$.
As a goal of this chapter, we will provide upper bounds on the complexity of the networks
making up the data-driven ROMs, where we especially outline a comparison between the POD-DL-ROM and the newly
introduced DOD-DL-ROM.

In order to analyse any error occuring from a data-based approximation, we need to set up a probability space, in whose
sense a convergence under an increased data set can be viewed. This is to be regarded as a clarification of 
Assumption (\ref{assumption sample}).
\begin{definition}[Sampling] \label{def: sampling}
    Let $(\Pcal, \mathcal{B}(\Pcal))$ be a Borel-measureable parameter space and
    $(\Omega_\prob, \Fcal, \prob)$ some underlying probability space. Then we refer to an 
    \emph{independent sampling of size $N$}, if for each $i \in  \{1, \dots, N\}$ the
    random variable mapping attaining a sample $(\mu_i, \nu_i)$
    \begin{equation*}
        X_i:(\Omega_\prob, \Fcal, \prob) \to (\Pcal, \mathcal{B}(\Pcal))
    \end{equation*}
    is measurable, its push forward $\prob^{X_i}$ admits the same distribution $\prob^X$ , i.e.
    \begin{equation*}
        \prob[X_i \in A] = \prob^{X_i}[A] = \prob^X[A],
    \end{equation*}
    and the samples are independent. We further name it \emph{iid sampling}, if the distribution $\prob^X$
    is the uniform distribution
    on $\Pcal$.
\end{definition}

With this in mind, we can view sampling as the realization of a random variable, that provides an element of the sample
space. Therefore any mention of a probability space or a probability measure $\prob$ refers to the samplying procedure in the probability space
as given in Definition (\ref{def: sampling}). 

The upcoming proposition is a necessary result for the following chapters, as it provides us with a practical
bound on the needed number of samples.

\begin{proposition} \label{prop: order of convergence for sample}
    Let $f \in L^2(\Theta \times \Theta' \times [0, T])$. 
    Under Assumption (\ref{assumption sample}), it is true, that
    \begin{equation*}
        \E\left[ \bigg\vert \int_{\Theta \times \Theta' \times [0, T]} f(\mu, \nu, t) d(\mu, \nu, t) 
        - \frac{\vert \Theta \times \Theta' \times [0, T] \vert}{N_{data}} 
        \sum_{i = 1}^{N_{s_1}} \sum_{j = 1}^{N_{s_2}} \sum_{k = 1}^{N_t} f(\mu_i, \nu_j, k \Delta t) \bigg\vert \right] 
        \! \leq \mathcal{O}(N_s^{-1/2} + N_t^{-1}).
    \end{equation*}
\end{proposition}
\notdone{prove this}




\section{Error Decomposition Formulas} \label{section: error decomposition}




To adequately start analysing the ROMs in question, one needs to recognize the different aspects of "error" 
inherent to their deployment. Namely, we can subdivide these seperate instances of inaccuracy occuring in the process
of both, the POD-based data-driven ROM and the DOD-based data-driven ROM, into an error caused by the imperfect
recollection of the parameter spaces $\Theta$ and $\Theta'$, an error caused by the imperfect 
projection given by the DOD or POD respectively and an error of imperfect attribution of latent dynamics
on that projection space by a neural network. This seperation can be achieved in a similar manner for both types of algorithms.

Before we start with the formulation and proof of these decompositions, we want to predefine some elements of the
corresponding theorems, to make them less cumbersome to read.

Firstly let in general $n, N', N, N_A$ and $N_h$ be the fixed natural number hierarchy from 
Assumption (\ref{assumption: dimension hierarchy}). 
Let the approximate POD-based parameter-to-solution map be defined as follows
\begin{equation*}
    \Gcal_\text{POD+DNN}: \Theta \times \Theta' \times [0, T] \to \R^{N_h}; \quad (\mu, \nu, t) \mapsto 
    \A_\text{P} \cdot \hat{q}_\text{P}(\mu, \nu, t),
\end{equation*}
where we let $\A_\text{P} \in \R^{N_h \times N}$ be a POD computed by Algorithm (\ref{algo: pod}) for the dimension $N$ 
with the mass matrix $\G$,
and
\begin{equation*}
    \hat{q}_\text{P} := R_\rho(\Phi^\text{P}_N(\theta_\text{PDNN})) \in \R^N
\end{equation*}
be the realization of some neural network $\Phi^\text{P}_N$ with input dimension $p + q + 1$ and output dimension $N$,
and an already fixed weight $\theta_\text{PDNN} \in \Theta(\Phi^\text{P}_N)$. And let similarily the approximate
DOD-based parameter-to-solution map be defined as follows
\begin{equation*}
    \Gcal_\text{DOD+DNN}: \Theta \times \Theta' \times [0, T] \to \R^{N_h}; \quad (\mu, \nu, t) \mapsto 
    \V_{\mu, t} \cdot \hat{q}_\text{D}(\mu, \nu, t),
\end{equation*}
where we let 
\begin{equation*}
    \V_{\mu, t} := \A \cdot \tilde{\V}_{\mu, t}(\theta_\text{DOD}) \in \R^{N_h \times N'},
\end{equation*}
for an a priori fixed parameter $\theta_\text{DOD} \in \Theta(\Phi_{\tilde{\V}})$ according to Definition 
(\ref{def: DOD}) and the used POD matrix $\A \in \R^{N_h \times N_A}$, and
\begin{equation*}
    \hat{q}_\text{D} := R_\rho(\Phi_{N'}(\theta_\text{DDNN})) \in \R^{N'}
\end{equation*}
be the realization of some neural network $\Phi^\text{D}_{N'}$ with input dimension $p + q + 1$ and output dimension $N'$,
and an already fixed weight $\theta_\text{DDNN} \in \Theta(\Phi^\text{D}_{N'})$.

Moreover we assume for this section to have after reordering and with $N_\text{data} = N_{s_1} N_{s_2} N_t$
\begin{equation*}
    \left\{ (\mu_j, \nu_j, t_j), u_j \right\}_{j=1}^{N_\text{data}} :=
    \left\{(\mu_i, \nu_j, k \Delta t), u_h^{\mu_i, \nu_j, k \Delta t}\right\}_{(i, j, k)=(1, 1, 1)}^{(N_{s_1}, N_{s_2}, N_t)}
    \subset \left( \Pcal_1 \times \Pcal_2 \times [0, T] \right) \times \R^{N_h}
\end{equation*}
a set of training data conforming to the label and feature notation in Definition (\ref{def: feature and label space}),
where the labels have been sampled in accordance to Assumption (\ref{assumption sample}). 

For both algorithms, the relative error can be given as
\begin{equation} \label{eq: relative error}
    \Ecal_\text{R} := \left( \int_{\Theta \times \Theta' \times [0, T]} \frac{\norm{\uht - \uhthat}}{\norm{\uht}} 
    d(\mu, \nu, t)\right)^{1/2},
\end{equation}
where $\uhthat \in \R^{N_h}$ can be given either with regard to the approximate parameter-to-solution map
of the DOD+DNN or the POD+DNN.


\begin{theorem}[POD+DNN Decomposition] \label{theo: relative error POD+DNN}
    Define the preliminaries such as in the introduction to this Section (\ref{section: error decomposition}). 
    Then under the Assumptions (\ref{assumption sample}) and 
    (\ref{assumption parameter-to-solution map}), we have for the relative error of the POD+DNN
    
    \begin{equation*}
        \mathcal{E}_R^\text{POD+DNN} \leq \mathcal{E}_\text{S} + \mathcal{E}_{\text{POD}} + \mathcal{E}_{\text{NN}},
    \end{equation*}
    where 
    
    \begin{enumerate}
        \item The sampling error 
        \[
            \mathcal{E}_\text{S} = \mathcal{E}_\text{S}\left(\mathcal{G}, (\mu_i, \nu_i, t_i)_{i = 1}^{N_{\text{data}}}, N\right)
        \]
        satisfies $\mathcal{E}_\text{S} \to 0$ as $N_s, N_t \to \infty$, with expectation 
        \[
            \E[\mathcal{E}_\text{S}] = \mathcal{O}(N_s^{-1/4} + N_t^{-1}).
        \]
    
        \item The POD projection error 
        \[
            \mathcal{E}_{\text{POD}} = \mathcal{E}_{\text{POD}}
            \left(\mathcal{G}, (\mu_i, \nu_i, t_i)_{i = 1}^{N_{\text{data}}}, N\right)
        \]
        satisfies $\mathcal{E}_{\text{POD}} \to \mathcal{E}_{\text{POD}, \infty}$ almost surely as 
        $N_s, N_t \to \infty$, where
        \[
            \mathcal{E}_{\text{POD}, \infty} = \mathcal{E}_{\text{POD}, \infty}(\mathcal{G}, N)
        \]
        is independent of the data snapshots and proportional to the linear KnW of the manifold 
        $\Scal = \{ \uht \, \vert \, (\mu, \nu, t) \in \Theta' \times [0, T] \}$.
    
        \item The neural network approximation error
        \[
            \mathcal{E}_{\text{NN}} = \mathcal{E}_{\text{NN}}\left(\mathcal{G}, N, \theta_{PDNN}, \Phi_N^\text{P}\right)
        \]
        is arbitrarily small depending on the size and choice of weights of the neural network $\Phi^\text{P}_N$.
    \end{enumerate} 

\end{theorem}
\begin{proof}
    Let 
    \begin{equation*}
        \norm{\cdot}_{L^2_\omega} := \left(\int_{\Theta \times \Theta \times [0, T]} \frac{\norm{\cdot}^2}{\norm{\uht}^2} 
        d(\mu, \nu, t) \right)^{1/2},
    \end{equation*}
    define a map $\norm{\cdot}_{L^2_\omega}: L^2(\Theta \times \Theta' \times [0, T]; \R^{N_h}) \to \R^+$. Then this is
    a norm. We coarsely proof this, by remarking the inheritance of the triangle inequality from the $\norm{\cdot}$ under 
    the usage of the Minkowski inequality, i.e. for arbitrary $v, w \in L^2(\Theta \times \Theta' \times [0, T]; \R^{N_h})$
    \begin{align*}
        \norm{v + w}_{L^2_\omega}^2 &= \int_{\Theta \times \Theta' \times [0, T]} 
        \frac{\norm{v(\mu, \nu, t) + w(\mu, \nu, t)}^2}{\norm{\uht}^2}
        d(\mu, \nu, t) \\
        &\leq  \int_{\Theta \times \Theta' \times [0, T]} 
        \frac{\norm{v(\mu, \nu, t)}^2 + \norm{w(\mu, \nu, t)}^2}{\norm{\uht}^2}
        d(\mu, \nu, t) = \norm{v}_{L^2_\omega}^2 + \norm{w}_{L^2_\omega}^2,
    \end{align*}
    and the absolut homogeneity from simple qualities of the integral. Finally let
    $\norm{v}_{L^2_\omega} = 0$ for some $v \in L^2(\Theta \times \Theta' \times [0, T])$, then assuming $v \neq 0$ 
    yields with Assumption (\ref{assumption parameter-to-solution map})
    \begin{equation*}
        \norm{v}_{L^2_\omega}^2 \geq  m^{-1} \int_{\Theta \times \Theta' \times [0, T]} \norm{v(\mu, \nu, t)} d(\mu, \nu, t) 
        > 0,
    \end{equation*}
    hence reaching a contradiction.
    With the knowledge of $\norm{ \cdot }_{L^2_\omega}$ being a norm, we know it must be subject to the 
    triangle inequality, thus
    \begin{align}
        \mathcal{E}_\text{R} &= \left( \int_{\Theta \times \Theta' \times [0, T]} 
        \frac{\norm{\uht - \A_\text{P} \hat{q}_\text{P}(\mu, \nu, t)}}{\norm{\uht}} d(\mu, \nu, t) \right)^{1/2} \\ 
        \label{eq: decomposed errors poddnn}
        &\leq \norm{\uht - \A_\text{P} \A_\text{P}^T \G \uht}_{L_\omega^2} + \norm{\A_\text{P} \A_\text{P}^T \G \uht - 
        \A_\text{P} \hat{q}_\text{P}(\mu, \nu, t)}_{L^2_\omega}.
    \end{align}
    Let $q_N(\mu, \nu, t) := \A_\text{P}^T \G \uht \in \R^N$ be a reduced solution according to the POD basis. 
    Then the natural definition of the neural network approximation error is
    \begin{equation}
        \mathcal{E}_{\text{NN}} = \left( \int_{\Theta \times \Theta' \times [0, T]} \frac{\norm{\A_\text{P} 
        q_N(\mu, \nu, t) - 
        \A_\text{P} \hat{q}_\text{P}(\mu, \nu, t)}^2}{\norm{\uht}^2} d(\mu, \nu, t) \right)^{1/2}.
    \end{equation}
    This marks the second part in (\ref{eq: decomposed errors poddnn}). This is only dependent on the ability of 
    the neural network to identifying the best choice on a given linear sub-manifold of the solution manifold. 
    The first part is concerned with the error generated from the best \emph{possible} approximation of such a 
    linear sub-manifold.
    
    For this, we can bound the norm $\norm{ \cdot }_{L^2_\omega} \leq m^{-1} \norm{ \cdot }$, 
    where $m > 0$ is the constant stemming from Assumption (\ref{assumption parameter-to-solution map}). 
    Further we remind ourselves of the definition of the discrete correlation matrix 
    $K = \vert \Theta \times \Theta \times \mathcal{T} \vert N_{data}^{-1} U \G U^T$ and its eigenvalues 
    $\sigma_k^2$ in (\ref{def: POD}). Using $\sqrt{a + b} \leq \sqrt{a} + \sqrt{b}$ for all $a,b \geq 0$, we get
    \begin{align*}
        &\norm{\uht - \A_\text{P} \A_\text{P}^T \G \uht}_{L^2_\omega} \\
        &\leq m^{-1} \left( \int_{\Theta \times \Theta' \times [0, T]} \norm{\uht - \A_\text{P} \A_\text{P}^T \G \uht}^2 
        d(\mu, \nu, t) \right)^{1/2} \\
        &\leq m^{-1} \left( \int_{\Theta \times \Theta' \times [0, T]} \norm{\uht - \A_\text{P} \A_\text{P}^T \G \uht}^2 
        d(\mu, \nu, t) - \sum_{k > N} \sigma_k^2  + \sum_{k > N} \sigma_k^2 \right)^{1/2} \\
        &\leq m^{-1} \left( \, \bigg\vert \, \int_{\Theta \times \Theta' \times [0, T]} \norm{\uht - 
        \A_\text{P} \A_\text{P}^T \G \uht}^2 d(\mu, \nu, t) - \sum_{k > N} \sigma_k^2 \, \bigg\vert + \sum_{k > N} \sigma_k^2 \right)^{1/2} \\
        &\leq \underbrace{m^{-1} \, \bigg\vert \, \int_{\Theta \times \Theta' \times [0, T]} \norm{\uht - 
        \A_\text{P} \A_\text{P}^T \G \uht}^2 d(\mu, \nu, t) - \sum_{k > N} \sigma_k^2 \, \bigg\vert^{1/2}}_{
            =: \mathcal{E}_\text{S}} + \underbrace{m^{-1} \sqrt{\sum_{k > N} \sigma_k^2}}_{=: \mathcal{E}_{\text{POD}}}. \\
    \end{align*}
    This concludes the estimate up to the facts about the convergence of the errors. Recalling
    Proposition (\ref{prop: eigenvalues and optimal projection under sampling}), 
    we can rewrite the sampling error as follows
    \begin{equation*}
        \mathcal{E}_\text{S} = m^{-1} 
        \bigg| \int_{\Theta \times \Theta' \times [0, T]} 
        \norm{\uht - \A_\text{P} \A_\text{P}^T \G \uht}^2 \, d(\mu, \nu, t) \\
        - \frac{|\Theta \times \Theta' \times [0, T]|}{N_{data}} 
        \sum_{j=1}^{N_{data}} \norm{u_j - \A_\text{P} \A_\text{P}^T \G u_j}^2 \bigg|^{1/2}
    \end{equation*}
    Furthermore, we define, using the compactness postulated in Assumption (\ref{assumption sample}) and 
    the boundedness in Assumption (\ref{assumption parameter-to-solution map}), the $L^2$-error map via
    \begin{equation*}
        f(\mu, \nu, t) := \norm{\uht - \A_\text{P} \A_\text{P}^T \G \uht}^2 \leq M^2 \opnorm{I_{N_h} - \A_\text{P} \A_\text{P}^T \G}^2 = 
        M^2 \opnorm{I_{N_A} - \A_\text{P}} < \infty.
    \end{equation*}
    Thus with $f \in L^2(\Theta \times \Theta' \times [0, T])$, we know that by Proposition 
    (\ref{prop: order of convergence for sample}), $\E[\mathcal{E}_\text{S}] \leq \mathcal{O}(N_s^{-1/4} + N_t^{-1})$.
    
    By the strong law of large numbers, we can thus conclude, that $\mathcal{E}_\text{S} \to 0$ almost surely, 
    as $N_t, N_s \to \infty$, and further by Proposition (\ref{prop: existence and convergence to sigma infty}), 
    we have that $\mathcal{E}_{POD} \to \mathcal{E}_{POD, \infty}$ as $N_t, N_s \to \infty$, where
    \begin{equation} \label{eq: pod eigenvalues cts case}
        \mathcal{E}_{POD, \infty} = m^{-1} \sqrt{\sum_{k > N} \sigma_{k, \infty}^2}.
    \end{equation}
\end{proof}

Even optimistically, the upper error bound decomposition of Theorem (\ref{theo: relative error POD+DNN}) suggest
for an ever-increasing number of data points $N_s, N_t \to \infty$, and an ever-increasing complexity
for the Deep Neural Network underlying the realization of $\hat{q}_\text{P}$, an error of order of the 
Kolmogorov $N$-width of the solution manifold, since the DNN cannot overcome the unavoidable loss to
even the perfect POD choice. Intuitively, the specific proportionality constant will be larger if the 
quotient between "small" regions of the solution space to the "large" regions of the solution space is big.

This observation can be reformulated into the following theorem.

\begin{theorem} \label{theo: lower bound pod+dnn}
    Under the assumptions of Theorem (\ref{theo: relative error POD+DNN}), we have that 
    \begin{equation*}
        \Ecal_\text{R}^\text{POD+DNN} \geq \frac{m}{M} \Ecal_{\text{POD}, \infty},
    \end{equation*}
    where $\Ecal_{\text{POD}, \infty} := m^{-1} \sqrt{\sum_{k > N} \sigma_{k, \infty}^2}$.
\end{theorem}
\begin{proof}
    We can simply start from the definition and argue, that under Proposition (\ref{prop: eigenvalues and optimal projection under sampling})
    the projection coefficients are chosen at best sub-optimal, i.e.
    \begin{equation*}
        \Ecal_\text{R} = \int_{\Theta \times \Theta' \times [0, T]} 
        \frac{\norm{\uht - \A_\text{P} \hat{q}_N(\mu, \nu, t)}}{\norm{\uht}} d(\mu, \nu, t)
        \geq
        \int_{\Theta \times \Theta' \times [0, T]} 
        \frac{\norm{\uht - \A_\text{P} \A_\text{P}^T \G \uht}}{\norm{\uht}} d(\mu, \nu, t),
    \end{equation*}
    where $\A_\text{P} \in \R^{N_h \times N}$ is the POD matrix as given by Definition (\ref{def: POD}). Then
    starting from the other direction for a POD matrix $\A_\infty$ under the continuous correlation matrix, we
    get under the definition of the proof of Theorem (\ref{theo: relative error POD+DNN})
    \begin{align}
        (\Ecal_{\text{POD}, \infty})^2 &= m^{-2} \sum_{k > N} \sigma_{k,\infty}^2 \\
        &= m^{-2} \int_{\Theta \times \Theta' \times [0, T]} \norm{\uht - \A_\infty \A_\infty^T \G \uht} d(\mu, \nu, t)
        \label{align: lower bound 1} \\
        &\leq m^{-2} \int_{\Theta \times \Theta' \times [0, T]} \norm{\uht - \A_\text{P} \A_\text{P}^T \G \uht} d(\mu, \nu, t)
        \label{align: lower bound 2} \\
        &= m^{-2} \int_{\Theta \times \Theta' \times [0, T]} \frac{\norm{\uht - \A_\text{P} \A_\text{P}^T \G \uht}}{\norm{\uht}}
        \norm{\uht} d(\mu, \nu, t) \\
        &\leq \frac{M^2}{m^2} (\Ecal_\text{R})^2, \label{align: lower bound 3}
    \end{align}
    where we use the Proposition (\ref{prop: eigenvalues and optimal projection under sampling}) for the
    continuous POD matrix in Step (\ref{align: lower bound 1}), the optimality of the same proposition 
    for Step (\ref{align: lower bound 2}) and the first equation of this proof in 
    Step (\ref{align: lower bound 3}) with the rather obvious use of the bound on $\norm{\uht}$ from Assumption
    (\ref{assumption parameter-to-solution map}).
\end{proof}


We have now discussed some preliminary relative error results for the POD+DNN neural networks for given near-optimal
neural network weights. As a next step, one can do the same for the DOD+DNN neural networks. First, we need
an auxilary statement, that precidents the similarity of definition of the DOD projection error to the POD projection
error, since the former depends on architecture. 

The key idea here being, that one can employ Horniks Theorem to ensure existence of a neural network architecture and corresponding neural 
network of such architecture, such that the objective functional 
\begin{equation*}
    (\mu, t, \V) \mapsto \int_{\Theta'} \norm{\uht - \V_{\mu, t} \V_{\mu, t}^T \G \uht} d (\nu),
\end{equation*}
is close to the "perfect one". Under triangle equality for the distance to the "perfect one", this can be seperated
into the mistake, which is arbitrarily small, and the total size of the perfect objective functional, that is nothing
but the Kolmogorov $N'$-width under fixed $(\mu, t)$. And that KnW decays exponentially fast by Assumption 
(\ref{assumption KnW decay}), in contrast to the general one. 

\begin{remark} \label{remark: regularity for s}
    Generally we can \emph{not} employ Yarotsky Theorem (\ref{theo: yarotsky}) here, since even though 
    the objective
    functional, as described above, can be shown to be Lipschitz continous under our current 
    assumptions, as given above, its minimizer with regard to the matrix entry does not 
    necessarily inherit any regularity and Yarotsky
    does require Sobolev-smoothness. However, we will set up the proof for the following lemma in such a way,
    that Yarotsky can be directly interchanged, if a higher regularity for the optimal objective functional is
    imposed.
\end{remark}

Before we start, we would like to define the aforementioned optimal objective functional, since further assumptions
on the PDE can expose the needed further regularity.

The source for the following proposition and its framework is generally \cite{franco2024}.

\begin{definition}[Optimal DOD Functional] \label{def: DOD functional}
    Let $N', N_A, N_h \in \N$ be fixed by Assumption (\ref{assumption: dimension hierarchy}). 
    Assuming we have $\mu \in \Theta, t \in [0, T]$ and $\V = \A \tilde{\V} \in [0, 1]^{N_h \times N'}$, with
    $\A \in [0, 1]^{N_h \times N_A}$ being $\G$ orthonormal, and $\tilde{\V} \in [0, 1]^{N_A \times N'}$, then 
    we define $J: \Theta \times [0, T] \times [0, 1]^{N_h \times N'} \to \R$ via
    \begin{equation*}
        J(\mu, t, \V) := \int_{\Theta'} \norm{\uht - \V \V^T \G \uht} d(\nu).
    \end{equation*}
    Then we call a Borel measurable map $s: \Theta \times [0, T] \to [0, 1]^{N_A \times N'}$ 
    producing w.l.o.g. orthonormal
    matrices given by
    \begin{equation*}
        s: (\mu, t) \mapsto \underset{{\tilde{\V} \in [0, 1]^{N_A \times N'}}}{\text{argmin}}  
        J(\mu, t, \A \tilde{\V}),
    \end{equation*}
    the \emph{optimal DOD functional}, if it is uniquely defined.
\end{definition}

Under the already given assumptions, we can pose its first existence result.
\begin{lemma} \label{lemma: DOD functional is lipschitz}
    Let $J: \Theta \times [0, T] \times [0, 1]^{N_h \times N'} \to \R$ be given as in 
    Definition (\ref{def: DOD functional}),
    then under Assumption (\ref{assumption parameter-to-solution map}), we have that $J$ is Lipschitz,
    and the optimal DOD functional $s$ exists.
\end{lemma}
\begin{proof}
    Let $\mu, \mu' \in \Theta$, $t, t' \in [0, T]$ 
    and $\V = \A \tilde{\V}, \V' = \A \tilde{\V}' \in \R^{N_h \times N'}$ such that for some $\delta > 0$
    \begin{equation*}
        \euknorm{\mu - \mu'} + \euknorm{t - t'} + \opnorm{\V - \V'} < \delta.
    \end{equation*}
    For the following calculation, we compose the preliminary identity for 
    $u := \uhtprered$ and $u' := u_\text{pre-red}^{\mu', \nu, t'}$,
    with $u_\text{pre-red}^{\mu, \nu, t} := \A^T \G \uht \in \R^{N_A}$ for each $\mu \in \Theta, \nu \in \Theta'$
    and $t \in [0, T]$, 
    and the projectors $P := \tilde{\V} \tilde{\V}^T$, $P' := \tilde{\V}' (\tilde{\V}')^T$,
    \begin{equation} \label{eq: optimal dod help equation}
        \begin{aligned}
            (\I_{N_A} - P)u - (\I_{N_A} - P')u' &= u - Pu - u' + P'u' \\
            &= u - Pu - u' + P'u' + Pu' - Pu' \\
            &= u - u' - Pu + Pu' + P'u' - Pu' \\
            &= (\I_{N_A} - P)(u - u') + (P' - P)u'.
        \end{aligned}
    \end{equation}
    Then we can calculate with, the reverse triangle inequality, and our auxiliary 
    Equation (\ref{eq: optimal dod help equation}),
    \begin{align*}
        \euknorm{J(\mu, t, \V) - J(\mu', t', \V')} &= \int_{\Theta'} \euknorm{\norm{\uht - \V \V^T \G \uht} - 
        \norm{u_h^{\mu', \nu, t} - \V' (\V')^T \G u_h^{\mu', \nu, t}}} d(\nu) \\
        &= \int_{\Theta'} \left\vert \, \euknorm{u_\text{pre-red}^{\mu, \nu, t} - \tilde{\V} 
        \tilde{\V}^T u_\text{pre-red}^{\mu, \nu, t}} - 
        \euknorm{u_\text{pre-red}^{\mu', \nu, t'} - \tilde{\V}' (\tilde{\V}')^T 
        u_\text{pre-red}^{\mu', \nu, t'}} \, \right\vert d(\nu) \\
        &= \int_{\Theta'} \left\vert \, \left(u_\text{pre-red}^{\mu, \nu, t} - \tilde{\V} 
        \tilde{\V}^T u_\text{pre-red}^{\mu, \nu, t}\right) - 
        \left(u_\text{pre-red}^{\mu', \nu, t'} - \tilde{\V}' (\tilde{\V}')^T 
        u_\text{pre-red}^{\mu', \nu, t'}\right) \, \right\vert d(\nu) \\
        &\leq \int_{\Theta'} \left\vert (\I_{N_A} - \tilde{\V} \tilde{\V}^T)(\uhtprered - 
        u_\text{pre-red}^{\mu', \nu, t'} )  \right\vert d(\nu) + M \vert \Theta' \vert \opnorm{\tilde{\V} \tilde{\V}^T - 
        \tilde{\V}' (\tilde{\V}')^T}   \\
        &< \delta \vert \Theta' \vert (L  + 2  M ),
    \end{align*}
    and $J$ is thus Lipschitz continous.
    \notdone{include existence of s}
\end{proof}



 

\begin{proposition}[Optimal Deep Orthogonal Decomposition] \label{prop: DOD Existence}
    For each geometric parameter $\mu \in \Theta$, let $\Scal_{\mu, t} :=
    \{ \uht \}_{\nu \in \Theta'} \subset \mathcal{G}(\Theta \times \Theta' \times [0, T])$ be the 
    $(\mu, t)$-invariant solution submanifold. 
    Then, under the Assumptions (\ref{assumption sample}) and (\ref{assumption parameter-to-solution map}), 
    for every $\varepsilon > 0$ there are weights $\theta_\text{DOD}^* \in \Theta(\Phi_{\tilde{\V}})$ 
    for a DOD of output dimension $N_h \times N'$ accoring to Definition (\ref{def: DOD}) such that
    \begin{equation*}
        \int_{\Theta \times \Theta' \times [0, T]} 
        \norm{\uht - \V_{\mu, t}(\theta_\text{DOD}^*) \V_{\mu, t}(\theta_\text{DOD}^*)^T \G \uht} d(\mu, \nu, t)
        < \varepsilon + \int_{\Theta \times [0, T]} d_{N'}(\mathcal{S}_{\mu, t}) d(\mu, t).
    \end{equation*}
\end{proposition}
\begin{proof}
    Let $N', N_A, N_h \in \N$ be fixed by Assumption (\ref{assumption: dimension hierarchy}). 
    We set $J$ and $s$ accoring to Definiton (\ref{def: DOD functional}), whereas the map
    \begin{equation*}
        s: (\mu, t) \mapsto \underset{{\tilde{\V} \in [0, 1]^{N_A \times N'}}}{\text{argmin}}  J(\mu, t, \A \tilde{\V}).
    \end{equation*}
    is bounded (given by continuity of J and compactness of its image) and therefore 
    $s \in L^1(\Theta \times [0, T]; \R^{N_A \times N'})$, i.e.
    \begin{equation*}
        \norm{s}_{L^1(\Theta \times [0, T]; \R^{N_A \times N'})} := 
        \int_{\Theta \times [0, T]} \opnorm{s(\mu, t)} \, d(\mu, t)  < \infty.
    \end{equation*}
    Define $[s]_{ij}$ as the canoncially lifted
    entry in position $(i, j) \in \{1, \dots, N_A\} \times \{1, \dots, N'\}$.

    This means that by Hornik Theorem (\ref{theo: hornik theorem}), there is a neural network architecture 
    $\Acal^{[\tilde{\V}_0]_{ij}}$ with  $(p+1)$-dimensional input and one-dimensional ouput, 
    and a neural network of that architecture $[\tilde{\V}_0]_{ij}$, satisfying
    \begin{equation*}
        \norm{[s]_{ij} - R_\rho([\tilde{\V}_0]_{ij})}_{L^1(\Theta \times [0, T]; \R)} \leq \varepsilon,
    \end{equation*}
    for each $(i, j) \in \{1, \dots, N_A\} \times \{1, \dots, N'\}$, and then using network parallelization 
    as in Definition (\ref{def: network parallelization}), we can set the optimal DOD neural network as
    \begin{equation*}
        \tilde{\V}_0 := P\left(P([\tilde{\V}_0]_{1 1}, \dots, [\tilde{\V}_0]_{1 N'}), \dots, 
        P([\tilde{\V}_0]_{N_A  1}, \dots,
        [\tilde{\V}_0]_{N_A N'})\right),
    \end{equation*}
    since it attains the wanted accuracy
    \begin{equation*}
        \norm{s - R_\rho(\tilde{\V}_0)}_{L^1(\Theta \times [0, T]; \R^{N_A \times N'})} \leq \varepsilon.
    \end{equation*}
    To now ensure, that the provided realization of the DOD neural network are indeed in the pre-image of $J$, 
    we need to ensure, that there is a network, which indeed attains the accuracy, but is also entrywise $[0,1]$.
    For this, let $L = ((A_1, b_1), (A_2, b_2))$ be given via
    \begin{equation*}
        A_1 := \begin{bmatrix}
            \I_{\R^{N_A \times N'}} \\
            \I_{\R^{N_A \times N'}}
        \end{bmatrix},
        \quad
        b_1 := \begin{pmatrix}
            0_{\R^{N_A \times N'}} \\
            -\I_{\R^{N_A \times N'}}
        \end{pmatrix},
        \quad
        A_2 := \begin{bmatrix}[c|c]
            \I_{\R^{N_A \times N'}} & -\I_{\R^{N_A \times N'}} 
        \end{bmatrix},
        \quad
        b_2 := 0.
    \end{equation*}
    The realization of this network is equivalent to the entry-wise operator of $x \mapsto \rho(x)-\rho(x-1) \in [0, 1]$
    for $x \in \R$. Since this realization as a map is bounded by $1$, we can infer with the definition
    $\V_{\mu, t} := R_\rho(L(\theta^*_\text{DOD})) \circ (\A \cdot R_\rho(\tilde{\V}_0(\theta^*_\text{DOD}))) 
    (\mu, t)$ a new network with a final
    architecture $\Phi_{\tilde{\V}}$ and
    complexity only differing by constants, and a choice of neural network weights $\theta^*_\text{DOD} 
    \in \Theta(\Phi_{\tilde{\V}})$, such that
    \begin{equation*}
        \int_{\Theta \times [0, T]} \opnorm{\V_{\mu, t} - s(\mu, t)} d(\mu, t)
        = \int_{\Theta \times [0, T]} \underbrace{\opnorm{L}}_{\leq 1} \opnorm{R_\rho(\tilde{\V}_0)(\mu, t) - s(\mu, t)}
        d(\mu, t)
        < \varepsilon.
    \end{equation*}
    Finally, with the knowledge, that $\V_{\mu, t}$ is $\varepsilon$-close to $s(\mu, t)$ in average, and the fact
    that this, as a matrix, is indeed in the pre-image of $J$, we can write using the Lipschitz continuity of $J$
    \begin{multline} \label{eq: last step for optimal DOD}
        \int_{\Theta \times \Theta' \times [0, T]} 
        \norm{\uht - \V_{\mu, t} \V_{\mu, t}^T \G \uht} d(\mu, \nu, t) \\ 
        \leq \int_{\Theta \times [0, T]}
        \left\vert J(\mu, t, \V_{\mu, t}) - J(\mu, t, s(\mu, t)) \right\vert d(\mu, t) +
        \int_{\Theta \times [0, T]} J(\mu, t, s(\mu, t)) d(\mu, t) \\
        < \varepsilon + \int_{\Theta \times [0, T]} \left( \min\limits_{\V \in \R^{N_h \times N}} J(\mu, t, \V)\right) 
        d(\mu, t) \\
        < \varepsilon + \int_{\Theta \times [0, T]} \left( \min\limits_{\V \in \R^{N_h \times N}} 
        \sup\limits_{\nu \in \Theta'} \norm{\uht - \V \V^T \G \uht} \right) d(\mu, t).
    \end{multline}
\end{proof}

\begin{remark} \label{remark: discrete optimal DOD}
    We have structured the Proposition (\ref{prop: DOD Existence}) with the $L^1$-integral in mind, but we 
    may redefine it to $L^\infty$ and use the following definition of $J$ for the statement and proof
    \begin{equation*}
        J(\mu, t, \V) := \frac{\euknorm{\Theta'}}{N_{s_2}} \sum_{\nu \in \Pcal_2} \norm{\uht - \V \V^T \G \uht}.
    \end{equation*}
    Then the entire proof can be done in the same way, since the Lipschitz property of $J$ in Lemma
    (\ref{lemma: DOD functional is lipschitz}) remains true, 
    and culminates in the following version of Proposition (\ref{prop: DOD Existence}) for $\V_{\mu, t} :=
    \V_{\mu, t}(\theta^*_\text{DOD})$, when discarding the very last step in 
    Equation (\ref{eq: last step for optimal DOD})
    \begin{multline*}
        \sup\limits_{(\mu, t) \in \Theta \times [0, T]} \frac{\euknorm{\Theta'}}{N_{s_2}} \sum_{\nu \in \Pcal_2}
        \norm{\uht - \V_{\mu, t}
        \V_{\mu, t}^T \G \uht} \\
        < \varepsilon + \sup\limits_{(\mu, t) \in \Theta \times [0, T]} \frac{\euknorm{\Theta'}}{N_{s_2}} 
        \min\limits_{\V \in \R^{N_h \times N'}} 
        \sum_{\nu \in \Pcal_2} 
        \norm{\uht - \uht \V \V^T \G \uht}.
    \end{multline*}
\end{remark}


With that in mind, one can now build an analogon to the relative error on the
POD+DNN ROM (i.e. Theorem (\ref{theo: relative error POD+DNN})), since we now have a similar statement
as the one of Proposition (\ref{prop: eigenvalues and optimal projection under sampling}), in the form
of Remark (\ref{remark: discrete optimal DOD}). In fact, we can even leverage 
Proposition (\ref{prop: eigenvalues and optimal projection under sampling}) to define
\begin{equation} \label{eq: definition virtual eigenvalues}
    \sigma(\mu, t)_1^2 \geq \cdots \geq \sigma(\mu, t)_{n}^2 \quad \text{for any } n \leq r,
\end{equation}
as the eigenvalues of a hypothetic POD (see Definition (\ref{def: POD})) with the snapshot matrix
\begin{equation*}
    U := \left[u_h^{\mu, \nu_1, t} \big\vert \cdots \big\vert u_h^{\mu, \nu_{N_{s_2}}, t} \right]^T 
    \in \R^{N_h \times N_{s_{2}}},
\end{equation*}
for each $(\mu, t) \in \Theta \times [0, T]$.

Indeed, intuitively this makes sense, since the DOD is a neural network, and can achieve the optimal projector 
out of sheer luck, independent of the amount of training data provided.



\begin{theorem} \label{theo: relative error DOD+DNN}
    Define the preliminaries such as in the introduction to this Section (\ref{section: error decomposition}). 
    Further set $\V_{\mu, t} := \V_{\mu, t}(\theta^*_\text{DOD})$ for the choice of weights $\theta_\text{DOD}^* \in
    \Theta(\Phi_{\tilde{\V}})$ as in Proposition (\ref{prop: DOD Existence}).
    Then, under the Assumptions (\ref{assumption sample}) and 
    (\ref{assumption parameter-to-solution map}), we have for the relative error of the DOD+DNN
    
    \begin{equation}
        \mathcal{E}^\text{DOD+DNN}_R \leq \mathcal{E}_\text{S} + \mathcal{E}_\text{DOD} + \mathcal{E}_\text{NN},
    \end{equation}
    where 
    
    \begin{enumerate}
        \item The sampling error 
        \[
            \mathcal{E}_\text{S} = \mathcal{E}_\text{S}\left(\mathcal{G}, 
            (\mu_i, \nu_i, \tau_i)_{i = 1}^{N_{\text{data}}}, N'\right)
        \]
        satisfies $\mathcal{E}_\text{S} \to 0$ as $N_{s_2} \to \infty$, with expectation 
        \[
            \E[\mathcal{E}_\text{S}] = \mathcal{O}(N_{s_2}^{-1/4}).
        \]
    
        \item The DOD projection error 
        \[
            \mathcal{E}_{\text{DOD}} = \mathcal{E}_{\text{DOD}}\left(\mathcal{G}, 
            (\mu_i, \nu_i, \tau_i)_{i = 1}^{N_{\text{data}}}, N\right)
        \]
        satisfies $\mathcal{E}_{\text{DOD}} \to \mathcal{E}_{\text{DOD}, \infty}$ 
        almost surely as $N_s, N_t \to \infty$, where
        \[
            \mathcal{E}_{\text{DOD}, \infty} = \mathcal{E}_{\text{DOD}, \infty}(\mathcal{G}, N)
        \]
        is independent of the data snapshots and arbitrarily close, with regards to the size 
        of the DOD neural network, 
        to the order of the maximal linear KnW of the solution submanifold 
        $\mathcal{S}_{\mu, t} = \{ \uht \, \vert \, \nu \in \Theta' \}$.
    
        \item The neural network approximation error
        \[
            \mathcal{E}_{\text{NN}} = \mathcal{E}_{\text{NN}}\left(\mathcal{G}, N, \hat{q}\right)
        \]
        is arbitrarily small depending on the size and choice of weights of the neural network $\Phi_A$.
    \end{enumerate}

\end{theorem}
\begin{proof}
    
    In its core, this proof inherits its structure and arguments from the one of Theorem 
    (\ref{theo: relative error POD+DNN}), 
    we will therefore only shorten the parts in which they are interchangeable. Let 
    $\norm{\cdot}_{L^2_\omega}: L^2(\Theta \times \Theta' \times [0, T]; \R^{N_h}) \to \R^+$ be the norm as given before,
    and define
    \begin{equation*}
        q_N(\mu, \nu, t) := \V_{\mu, t}^T \G \uht \in \R^N, 
    \end{equation*}
    as the reduced solution regarding the DOD projection.
    Then we can apply the triangle inequality like before, to achieve 
    \begin{equation*}
        \Ecal_\text{R}^\text{DOD+DNN} \leq \norm{\uht - \V_{\mu, t} \V_{\mu, t}^T \G \uht}_{L^2_\omega} + 
        \norm{\V_{\mu, t} q_N(\mu, \nu, t) - \V_{\mu, t} \hat{q}_\text{D}(\mu, \nu, t)}_{L^2_\omega}.
    \end{equation*}
    In turn, from here we can already infer the error stemming from the latent dynamics identification under the 
    neural network realization $\hat{q}_\text{D}$, i.e.
    \begin{equation}
        \mathcal{E}_{NN} = \left( \int_{\Theta \times \Theta' \times [0, T]} 
        \frac{\norm{\V_{\mu, t} q_N(\mu, \nu, t) - \V_{\mu, t} \hat{q}_\text{D}(\mu, \nu, t)}^2}{\norm{\uht}^2} 
        d(\mu, \nu, t) \right)^{1/2}.
    \end{equation}
    
    Moreover, we can extrapolate a similar procedure of splitting the remaining term into the respective errors
    as in the proof of Theorem (\ref{theo: relative error POD+DNN}), 
    where we already assume for the fixed $N \in \N$, the "optimal" choice of weights $\theta_\text{DOD}^* \in
    \Theta(\Phi_{\tilde{\V}})$ for the underlying DOD neural network in accordance to Remark
    (\ref{remark: discrete optimal DOD}) and its subsequently defined virtual discrete 
    corrolation eigenvalues $\sigma(\mu, t)_k$
    in Equation (\ref{eq: definition virtual eigenvalues})
    \begin{align*}
        &\norm{\uht - \V_{\mu, t} \V_{\mu, t}^T \G \uht}_{L^2_\omega} \\
        &\leq m^{-1} \left( \int_{\Theta \times [0, T]} \int_{\Theta'} 
        \norm{\uht - \V_{\mu, t} \V_{\mu, t}^T \G \uht}^2 d(\nu) d(\mu, t) \right)^{1/2} \\
        &= m^{-1} \left( \int_{\Theta \times [0, T]} \int_{\Theta'} 
        \norm{\uht - \V_{\mu, t} \V_{\mu, t}^T \G \uht}^2 d(\nu) - \sum_{k > N'} \sigma(\mu, t)_k^2 +
        \sum_{k > N'} \sigma(\mu, t)_k^2 \, d(\mu, t) \right)^{1/2} \\
        &\leq \underbrace{m^{-1} \left\vert \, \int_{\Theta \times [0, T]} \int_{\Theta'} 
        \norm{\uht - \V_{\mu, t} \V_{\mu, t}^T 
        \G \uht}^2 d(\nu) -  \frac{\euknorm{\Theta'}}{N_{s_2}} \sum_{\nu \in \Pcal_2} \norm{\uht - \V_{\mu, t} \V_{\mu, t}^T \G \uht}^2 
        d(\mu, t) \, \right\vert^{1/2}}_{=: \Ecal_\text{S}} \\
        &+ 
        \underbrace{m^{-1} \euknorm{\Theta \times [0, T]}^{1/2}\varepsilon + m^{-1} \euknorm{\Theta \times [0, T]}^{1/2}
        \sup\limits_{(\mu, t) \in \Theta \times [0, T]} \sqrt{\sum_{k > N'} \sigma(\mu, t)_k^2}}_{=: \Ecal_\text{DOD}}.
    \end{align*}

    Deploying the same strategy as in the proof of Theorem (\ref{theo: relative error POD+DNN}), we
    can assert, that indeed $\Ecal_\text{S} \to 0$ almost surely by the strong law of large numbers, and
    that $\E[\Ecal_\text{S}] \leq \Ocal(N_{s_2}^{-1/4})$ by
    Proposition (\ref{prop: order of convergence for sample}) with $N_t = 1, N_{s_1} = 1$.

    Moreover, we have that 
    \begin{multline*}
        \Ecal_\text{DOD} = \frac{\euknorm{\Theta \times [0, T]}^{1/2}}{m} \left( \varepsilon + 
        \sup\limits_{(\mu, t) \in \Theta \times [0, T]} \sqrt{\sum_{k > N'} \sigma(\mu, t)_k^2}
        \right) \\
        \underset{\N_{s_2} \to \infty}{\longrightarrow} \frac{\euknorm{\Theta \times [0, T]}^{1/2}}{m} \left( \varepsilon + 
        \sup\limits_{(\mu, t) \in \Theta \times [0, T]} \sqrt{\sum_{k > N'} \sigma(\mu, t)_{k, \infty}^2}
        \right) =: \Ecal_{\text{DOD}, \infty},
    \end{multline*}
    which can be arbitrarily improved for a more complex choice of the inner DOD module architecture. This dependency
    on the complexity is suggested by Yarotsky Theorem (\ref{theo: yarotsky}), since the $L^\infty$ error 
    correlates with the 
    $\Hcal_h$ representation error by the compactness of the time-parameter space, that is given
    by Assumption (\ref{assumption sample}), but is only proveable under additional assumptions as argued in 
    Remark (\ref{remark: regularity for s}).
\end{proof}


To connect this line of thought to the prior discussion for the POD+DNN, we remind the reader of the definition 
of the virtual eigenvalues $\sigma(\mu, t)_k$ given in Equation (\ref{eq: definition virtual eigenvalues}).
Under Remark (\ref{remark: discrete optimal DOD}) we can infer the existence of an arbitrary
precise neural network realizing such a quasi-POD map, thus we can provide a similar result to Theorem 
(\ref{theo: lower bound pod+dnn}) for the DOD+DNN.







\section{Upper Complexity Bounds} \label{section: upper complexity bounds}





In this section, we will specify under the Assumption (\ref{assumption sample}), (\ref{assumption parameter-to-solution map}),
(\ref{assumption perfect embedding}) and (\ref{assumption: dimension hierarchy}) the upper complexity bounds for
our three main algorithms: the POD-DL-ROM, the DOD+DFNN and the DOD-DL-ROM, as introduced in
Section (\ref{section: data-driven ROMs}). 

However, to achieve actual complexity
bounds for the latter two ROMs, we will necessarily need to impose one further assumption for the regularity of the 
optimal objective
functional as discussed in Remark (\ref{remark: regularity for s}). 




\notdone{make discussion for regularity of parameter-to-solution map}



For the following discussion, we want to again set some preliminaries to shorten the theorems layout.
Firstly, we assume the setup of the previous Section (\ref{section: error decomposition}), where we recall
defining a sample to be used for all subsequent algorithms both POD-based and DOD-based neural networks,
respectively marked POD+DNN and DOD+DNN. In fact,
set up $N$ and $N'$ as follows
\begin{equation} \label{eq: definition of N, N'}
    N = \underset{n \in \N}{\operatorname{arg \, min}} \left( \sum_{k > n} \sigma_k^2 \leq \frac{m^2}{9} \varepsilon^2 \right),
    \quad N' = \underset{n \in \N}{\operatorname{arg \, min}} \left( \sup\limits_{(\mu, t) \in \Theta \times [0, T]} 
    \sum_{k > n} \sigma(\mu, t)_k^2 \leq \frac{m^2}{16 \euknorm{\Theta \times [0, T]}} 
    \varepsilon^2 \right)
\end{equation}
with $\varepsilon > 0$ yet unspecified, and $\sigma_k^2$ according to 
Definition (\ref{def: POD}) and $\sigma(\mu, t)_k^2$ to
Equation (\ref{eq: definition virtual eigenvalues}). Note, that under Assumption (\ref{assumption KnW decay}) the hierarchy
of Assumption (\ref{assumption: dimension hierarchy}) is still preserved, if the choice of $\varepsilon$ is small
enough and $N_\text{data}$ is large enough.

From this starting point, refine the definition of the approximate parameter-to-solution map for the POD+DNN into 
\begin{equation*}
    \Gcal_\text{POD-DL-ROM} : \Theta \times \Theta' \times [0, T] \to \R^{N_h}; \quad (\mu, \nu, t)
    \mapsto \A_\text{P} \cdot \left( \psi_N \circ \phi_n \right) (\mu, \nu, t),
\end{equation*}
the POD-DL-ROM under weights $\theta_\text{POD-DL-ROM} = (\theta_\text{D}, \theta_\text{DF})$ with 
\begin{equation*}
    \quad \phi_n := R_\rho(\Phi_n(\theta_\text{DF})): \R^{p+q+1} \to \R^n,
    \psi_N := R_\rho(\Psi_N(\theta_\text{D})): \R^n \to \R^{N},  
\end{equation*}
from the Subsection (\ref{subsec: pod-dl-rom}). Note, that $\A_\text{P} \neq \A$ in general, since like before
$\A_\text{P} \in \R^{N_h \times N}$
and $\A \in \R^{N_h \times N_A}$. Similarily, refine the approximate parameter-to-solution map for the DOD+DNN into
\begin{equation*}
    \Gcal_\text{DOD+DFNN} : \Theta \times \Theta' \times [0, T] \to \R^{N_h}; \quad (\mu, \nu, t)
    \mapsto \V_{\mu, t} \cdot \phi_{N'}(\mu, \nu, t),
\end{equation*}
the DOD+DFNN under weights $\theta_\text{DOD+DFNN} = (\theta_\text{DOD}, \theta_{\text{DF}_2})$ with
\begin{gather*}
    \V_{\mu, t} := \A \cdot R_\rho(\Phi_{\tilde{\V}}(\theta_\text{DOD}))(\mu, t) \in \R^{N_h \times N'},
    \qquad \text{for each } (\mu, t) \in \Theta \times [0, T] \text{, and } \\
    \phi_{N'} := R_\rho(\Phi_{N'}(\theta_{\text{DF}_2})): \R^{p+q+1} \to \R^{N'}
\end{gather*}
from the Subsection (\ref{subsec: DOD+DFNN}) and the DOD Definiton (\ref{def: DOD}), and
\begin{equation*}
    \Gcal_\text{DOD-DL-ROM} : \Theta \times \Theta' \times [0, T] \to \R^{N_h}; \quad (\mu, \nu, t)
    \mapsto \V_{\mu, t} \cdot (\psi_{N'} \circ \phi_n)(\mu, \nu, t),
\end{equation*}
the DOD-DL-ROM under weights $\theta_\text{DOD-DL-ROM} = (\theta_\text{DOD}, \theta_{\text{D}_2}, 
\theta_{\text{DF}_3})$ with
\begin{gather*}
    \V_{\mu, t} := \A \cdot R_\rho(\Phi_{\tilde{\V}}(\theta_\text{DOD}))(\mu, t) \in \R^{N_h \times N'},
    \qquad \text{for each } (\mu, t) \in \Theta \times [0, T] \text{, and } \\
    \psi_{N'} := R_\rho(\Psi_{N'}(\theta_{\text{D}_2})): \R^n \to \R^{N'}, \quad
    \phi_{n} := R_\rho(\Phi_n(\theta_{\text{DF}_3})): \R^{p+q+1} \to \R^n,
\end{gather*}
from the Subsection (\ref{subsec: dod-dl-rom}) and the DOD Definition (\ref{def: DOD}).

Furthermore, let $\psi_* \in C^s(\R^n; \R^N)$ and $\psi_*' \in C^{s'}(\R^N; \R^n)$ respectively be the maps attaining
the perfect embedding Assumption (\ref{assumption perfect embedding}), with $s \gg s' \geq 2$. Since $N > N'$ in
general, the perfect embedding is implied for $N'$ aswell; let $\chi_* \in C^s(\R^n; \R^{N'})$ 
and $\chi_*' \in C^{s'}(\R^{N'}; \R^n)$ be the maps attaining the equivalent for $N'$.

Finally set constants
\begin{equation*}
    C_1 := \sup\limits_{\vert \alpha \vert \leq s'} \sup\limits_{v \in \R^N} \vert \, D^\alpha \psi_*(v) \, 
    \vert, \quad 
    C_2 := \sup\limits_{\vert \alpha \vert \leq s} \sup\limits_{w \in \R^n} \vert \, D^\alpha \psi_*'(w) \, \vert,
\end{equation*} 
and $C_1', C_2'$ analogously for $\chi_*$ and $\chi_*'$. 




\begin{theorem} \label{theo: PAC bound for POD-DL-ROM}
    Let the preliminaries of the POD-DL-ROM be given as in the introduction of this Section 
    (\ref{section: upper complexity bounds}).
    Then there is $N_\text{data} = N_\text{data}(\delta, \varepsilon)$, 
    a constant $c = c(\Theta, \Theta', T, L, C_1, C_2, p, q, n, s, s')$ and 
    a choice of weights for the POD-DL-ROM approximate parameter-to solution map $\A_\text{P} \cdot \hat{q}_N :=
    \A_\text{P} \cdot (\psi_N \circ \phi_n): \R^{p+q+1} \to \R^{N_h}$ 
    such that 
    \begin{equation*}
        \prob\left[\Ecal^{POD-DL-ROM}_R < \varepsilon\right] > 1 - \delta,    
    \end{equation*}
    and the complexity is given by
    \begin{enumerate}
        \item[-] $L_{\Psi_N} = c \log(\varepsilon^{-1})$ layers,
        \item[-] $\omega_{\Psi_N} = cN \varepsilon^{-n/(s-1)} \log(\varepsilon^{-1})$ active weights,
    \end{enumerate}
    for the Decoder, and 
    \begin{enumerate}
        \item[-] $L_{\Phi_n} = c \log(\varepsilon^{-1})$ layers,
        \item[-] $\omega_{\Phi_n} = cn \varepsilon^{-(p+q+1)} \log(\varepsilon^{-1})$ active weights,
    \end{enumerate}
    for the Deep Feed-Forward Neural Network.
\end{theorem}
\begin{proof}
    We can directly bound $\Ecal_\text{S} = \Ecal_\text{S}(N_{s_1}, N_{s_2}, N_t)$ defined in the proof
    of Theorem (\ref{theo: relative error POD+DNN}) independently of $N$, under the 
    Assumption (\ref{assumption sample})
    by the Weak Law of Large Numbers, i.e. for all $1 > \delta > 0$ and $\varepsilon > 0$ there are 
    $N_{s_1}, N_{s_2}, N_t$, such that
    \begin{equation} \label{eq: bound on S error pod-dl-rom}
        \prob[\Ecal_\text{S}((N_{s_1}, N_{s_2}, N_t)) < \varepsilon / 3] > 1 - \delta.
    \end{equation}
    This result provides the "probability" part of the PAC (probably almost correct) bound, which we want to show. 
    In other words, provided this, the other two error types can be bounded without regard to the probability
    space underlying the sampling procedure.
    Bounding $\Ecal_{\text{NN}}$ and $\Ecal_{\text{POD}}$
    by $\varepsilon / 3$ respectively will then provide the wanted result.
    For this, by choosing $N$ as in Equation (\ref{eq: definition of N, N'}),
    we simply rewrite
    \begin{equation} \label{eq: bound on POD error pod-dl-rom}
        \Ecal_{\text{POD}} := m^{-1} \sqrt{\sum_{k > N} \sigma_k^2} \leq \frac{\varepsilon}{3}.
    \end{equation}
    To now bound the neural network approximation error, 
    we will need to simply rewrite the error in terms of a form accommodating a use of Gühring or Yarotsky Theorem;
    \begin{equation} \label{eq:nn error for pod-dl-rom}
        \begin{aligned}
            \Ecal_{\text{NN}} &= \left( \int_{\Theta \times \Theta' \times [0, T]} 
            \frac{\norm{\A_\text{P} q_N(\mu, \nu, t) - \A_\text{P} \hat{q}_N(\mu, \nu, t)}^2}{\norm{\uht}^2}
            \, d(\mu, \nu, t) \right)^{1/2} \\
            &\leq m^{-1} \left( \int_{\Theta \times \Theta' \times [0, T]} 
            \norm{\A_\text{P} q_N(\mu, \nu, t) - \A_\text{P} \hat{q}_N(\mu, \nu, t)}^2 \, d(\mu, \nu, t) \right)^{1/2} \\
            &\leq m^{-1} \left( \left| \Theta \times \Theta' \times [0, T] \right| 
            \sup_{(\mu, \nu, t) \in \Theta \times \Theta' \times [0, T]} 
            \norm{\A_\text{P} q_N(\mu, \nu, t) - \A_\text{P} \hat{q}_N(\mu, \nu, t)}^2 \right)^{1/2} \\
            &\leq m^{-1} \left( \left| \Theta \times \Theta' \times [0, T] \right| \cdot \opnorm{\A_\text{P}}^2 
            \sup_{(\mu, \nu, t) \in \Theta \times \Theta' \times [0, T]} 
            \norm{q_N(\mu, \nu, t) - \hat{q}_N(\mu, \nu, t)}^2 \right)^{1/2} \\
            &= m^{-1} \left| \Theta \times \Theta' \times [0, T] \right|^{1/2} 
            \sup_{(\mu, \nu, t) \in \Theta \times \Theta' \times [0, T]} 
            \norm{q_N(\mu, \nu, t) - \hat{q}_N(\mu, \nu, t)}.
        \end{aligned}
        \end{equation}
    whereas we split the latter error into both parts, since $\hat{q}_N = \psi_N \circ \phi_n$ per definition. 
    This means employing two types of error estimations
    in the matter of finding the respective optimal maps for the 
    feed forward estimation on the latent dynamics space via $\phi_*: \R^{(p+q+1)} \to \R^n$ 
    and the decoder lift via $\psi_*:\R^n \to \R^N$.
    \begin{enumerate}
        \item[-] Let $\mathcal{S}_N := \{q(\mu, \nu, t) = \A_\text{P}^T \G \uht \in \R^N \, \vert \, (\mu, \nu, t) 
        \in \Theta \times \Theta' \times [0, T] \}$, then the image of the optimal
        embedding $\mathcal{V}_n = \psi_*'(\mathcal{S}_N) \in \R^n$ is such that $\text{diam}(\mathcal{V}_n) 
        \leq C_1 L \cdot \text{diam}(\Theta \times \Theta' \times [0, T])$, given 
        the Assumption (\ref{assumption parameter-to-solution map}) and the preliminaries.
        In fact, using the definition of Sobolev spaces (\ref{def: sobolev space})
        \begin{align*}
            \sup\limits_{w, w' \in \mathcal{V}_n} \euknorm{w - w'} &= 
            \sup\limits_{v, v' \in \Scal_N} \euknorm{\psi_*'(v) - \psi_*'(v')}\\
            &\leq C_1 \sup\limits_{v, v' \in \Scal_N} \euknorm{v - v'} \\
            &\leq C_1 \sup\limits_{(\mu, \nu, t), (\mu', \nu', t') \in \Theta \times \Theta' \times [0, T]}
            \euknorm{\A^T \G (\uht - u_h^{\mu', \nu', t'})} \\
            &\leq C_1 \sup\limits_{(\mu, \nu, t), (\mu', \nu', t') \in \Theta \times \Theta' \times [0, T]}
            \opnorm{\A^T \G} \, \euknorm{\uht - u_h^{\mu', \nu', t'}} \\
            &\leq C_1 L \cdot \text{diam}(\Theta \times \Theta' \times [0, T]).
        \end{align*}
        
        This argument allows us to find that $\psi_* \in \Fcal_{s, n, \infty, C_1 L \cdot \text{diam}(\Theta \times
        \Theta' \times [0, T])}$, after unilateral prior transformation $x \mapsto x / \euknorm{x} \in (0, 1)^n$ of
        any input, 
        for the target space of the optimal function required by Theorem Gühring (\ref{theo: gühring}). Thus 
        there are neural network weights $\theta_\text{D}$ for $\psi_N = R_\rho(\Psi_N(\theta_\text{D})) 
        : \R^n \to \R^N$ fitting the network realization to that optimal function, such that
        \begin{equation*}
            \norm{\psi_* - \psi_N}_{W^{1, \infty}((0,1)^d)} \leq \frac{m}{6} 
            \euknorm{\Theta \times \Theta' \times [0, T]}^{-1/2} \varepsilon,
        \end{equation*}
        with $L_{\Psi_N} = c \log(\varepsilon^{-1})$ layers and $\omega_{\Psi_N} = c N 
        \varepsilon^{-n/(s-1)} \log(\varepsilon^{-1})$ active weights. For the sake of completeness, we mention that
        one would build $N$-many entry-wise neural networks in parallel and hence attain for each $\omega_{[\psi_N]_1}
        = c \varepsilon^{-n/(s-1)} \log(\varepsilon^{-1})$ active weights under Gühring, then employ the Lemma
        (\ref{lemma: network parallelization}) for the parallelization of these. Note, that this 
        does not influence the order of the layer number. A similar approach 
        was already argued in the proof of Proposition (\ref{prop: DOD Existence}).
        This statement is equivalent to
        \begin{gather} \label{eq: decoder error for pod-dl-rom}
            \sup\limits_{v \in \mathcal{V}_N} \norm{\psi_*(v) - \psi_N(v)} 
            < \frac{m}{6} \euknorm{\Theta \times \Theta' \times [0, T]}^{-1/2} \varepsilon, \\
            \operatorname{ess} \sup\limits_{v, v' \in \mathcal{V}_N} 
            \frac{\euknorm{(\psi_* - \psi_N)(v) - (\psi_* - \psi_N)(v)}}{\euknorm{v - v'}} 
            < \frac{m}{6} \euknorm{\Theta \times \Theta' \times [0, T]}^{-1/2} \varepsilon,
        \end{gather}
        by definition of the derivative and the fact, that Sobolev norms are the sum of both the equations above.
        Additionally this implies that, $\psi_N$ is Lipschitz with constant $C_3 := C_2 + m/6 \cdot 
        \euknorm{\Theta \times \Theta' \times [0, T]}^{-1/2}$, since
        \begin{equation*}
            \frac{\euknorm{(\psi_* - \psi_N)(v) - (\psi_* - \psi_N)(v)}}{\euknorm{v - v'}} \geq 
            \frac{\euknorm{\psi_N(v) - \psi_N(v')}}{\euknorm{v - v'}} - \frac{\euknorm{\psi_*(v) - \psi_*(v')}}{\euknorm{v - v'}}
        \end{equation*}
        by inverse triangle inequality for any $v, v' \in \mathcal{V}_N$ and subsequently employing the $C_2$-bound for
        $\psi_*$ from the preliminaries of this Section (\ref{section: upper complexity bounds}).


        \item[-] Now let $\phi_*(\mu, \nu, t) = \psi_*'(q(\mu, \nu, t))$ for 
        $(\mu, \nu, t) \in \Theta \times \Theta' \times [0, T]$, which is 
        Lipschitz continuous, with constant $L C_1$. In fact, we have for any pair $(\mu, \nu, t), (\mu', \nu', t')
        \in \Theta \times \Theta' \times [0, T]$,
        \begin{align*}
            \euknorm{\psi_*(\mu, \nu, t) - \psi_*(\mu', \nu', t')} &= \euknorm{\psi_*'(\A_\text{P}^T \G \uht) - 
            \psi_*'(\A_\text{P}^T \G u_h^{\mu', \nu', t'})} \\
            &\leq C_1 \euknorm{\A_\text{P}^T \G (\uht - u_h^{\mu', \nu', t'})} \\
            &\leq C_1 L \cdot \text{diam}(\Theta \times \Theta' \times [0, T]).
        \end{align*} 
        The Lipschitz continuity of this map then allows us to employ Theorem (\ref{theo: lipschitz implies sobolev}),
        which guarantees that $\phi_* \in W^{1, \infty}(\Scal_N)$, since $\text{diam}(\Scal_N) \leq L \cdot 
        \text{diam}(\Theta \times \Theta' \times [0, T])$ by the same arguments as in the $\mathcal{V}_n$-case.
        Hence, after entry-wise normalization, $\phi_* \in \Fcal_{1, p+q+1, \infty, \text{diam}(\Scal_N)}$.

        Now we employ Theorem Yarotski (\ref{theo: yarotsky}) to assert the existence of neural
        network weights $\theta_\text{DF}$ such that its realization $\phi_n := R_\rho(\Phi_n(\theta_\text{DF})): 
        \R^{p+q+1} \to \R^n$ attains
        \begin{equation} \label{eq: dfnn error for pod-dl-rom}
            \norm{\phi_* - \phi_n}_{W^{0, \infty}(\Theta \times \Theta' \times [0, T])}
            < \frac{m}{6 C_3}\euknorm{\Theta \times \Theta' \times [0, T]}^{-1/2} \varepsilon,
        \end{equation}
        with $L_{\Phi_n} = c \log(\varepsilon^{-1})$ layers and 
        $\omega_{\Phi_n} = c n \varepsilon^{-(p+q+1)} \log(\varepsilon^{-1})$ active weights, again with
        the argument of neural network parallelization.
    \end{enumerate}
    Starting from the last inequality in (\ref{eq:nn error for pod-dl-rom}) we can perform the split in a clean fashion
    using the triangle inequality and reminding ourselves of the Assumption (\ref{assumption perfect embedding}), which
    guarantees $q = \psi_*(\psi_*')$, 
    and thus gives with both neural network estimates in Equations (\ref{eq: decoder error for pod-dl-rom})
    and (\ref{eq: dfnn error for pod-dl-rom}) 
    \begin{align*}
        &\sup_{(\mu, \nu, t) \in \Theta \times \Theta' \times [0, T]} 
        \norm{q(\mu, \nu, t) - \hat{q}_N(\mu, \nu, t)} \\
        &\leq \sup\limits_{(\mu, \nu, t) \in \Theta \times \Theta' \times [0, T]} 
            ( \, \norm{\psi_*(\psi_*'(\mu, \nu, t)) - \psi(\phi_*(\mu, \nu, t))} + 
            \norm{\psi(\phi_*(\mu, \nu, t)) - \psi(\phi(\mu, \nu, t))} \, )  \\
        &\leq 
            \sup\limits_{v \in \mathcal{V}_N} \norm{\psi_*(v) - \psi(v)} + 
            C_3 \sup\limits_{(\mu, \nu, t) \in \Theta \times \Theta' \times [0, T]} 
            ( \, \norm{\psi(\phi_*(\mu, \nu, t)) - \psi(\phi(\mu, \nu, t))} \, )  \\
        &< \frac{m}{6}\euknorm{\Theta \times \Theta' \times [0, T]}^{-1/2} \varepsilon + 
        C_3 \frac{m}{6 C_3}\euknorm{\Theta \times \Theta' \times [0, T]}^{-1/2} \varepsilon = 
        \frac{\varepsilon}{3}\, m \,\euknorm{\Theta \times \Theta' \times [0, T]}^{-1/2},
    \end{align*}
    which cancels with $m^{-1} \left| \Theta \times \Theta' \times [0, T] \right|^{1/2}$ to get 
    \begin{equation} \label{eq: bound on NN error for pod-dl-rom}
        \Ecal_{\text{NN}} < \frac{\varepsilon}{3}.
    \end{equation}
    Now putting together all three estimates (\ref{eq: bound on S error pod-dl-rom}), 
    (\ref{eq: bound on POD error pod-dl-rom}) and (\ref{eq: bound on NN error for pod-dl-rom}) 
    with the reasoning of Theorem (\ref{theo: relative error POD+DNN}) finally gives
    \begin{equation*}
        \Ecal_R \leq \Ecal_{S} + \Ecal_{\text{POD}} + \Ecal_{\text{NN}} < 
        \frac{\varepsilon}{3} + \frac{\varepsilon}{3}
        + \frac{\varepsilon}{3} = \varepsilon,
    \end{equation*}
    with higher probability than $1 - \delta$.
\end{proof}


Building upon the discussion in the beginning of this Section (\ref{section: upper complexity bounds}), we
may explore the implication of a more regular parameter-to-solution map $\Gcal$. Namely, if that map
satisfies differentiability, then the optimal objective functional $s$ of Definiton (\ref{def: DOD functional})
inherits it, and subsequently provides the possibility of the use of Yarotsky Theorem in Proposition 
(\ref{prop: DOD Existence}) and by inclusion Remark (\ref{remark: discrete optimal DOD}).


\notdone{Include an additional assumption to provide regularity for the optimal dod functional}


\begin{theorem} \label{theo: PAC bound for DOD-DL-ROM}
    Let the preliminaries of the DOD-DL-ROM be given as in the introduction of this Section 
    (\ref{section: upper complexity bounds}).
    Then there is $N_\text{data} = N_\text{data}(\delta, \varepsilon)$, 
    a constant $c = c(\Theta, \Theta', T, L, C_1, C_2, p, q, n, s, s')$ and 
    a choice of weights for the DOD-DL-ROM approximate parameter-to solution map $(\mu, \nu, t) \mapsto 
    \V_{\mu, t} \cdot \hat{q}_{N'} :=
    \V_{\mu, t} \cdot (\psi_{N'} \circ \phi_n): \R^{p+q+1} \to \R^{N_h}$ 
    such that 
    \begin{equation*}
        \prob\left[\Ecal^{DOD-DL-ROM}_R < \varepsilon\right] > 1 - \delta,    
    \end{equation*}
    and the complexity is given by
    \begin{enumerate}
        \item[-] $L_{\Phi_{\tilde{\V}}} = c \log(\varepsilon^{-1})$ layers,
        \item[-] $\omega_{\Phi_{\tilde{\V}}} = cN_A N' \varepsilon^{-(p+1)} \log(\varepsilon^{-1})$ active weights,
    \end{enumerate}
    for the inner DOD module,
    \begin{enumerate}
        \item[-] $L_{\Psi_{N'}} = c \log(\varepsilon^{-1})$ layers,
        \item[-] $\omega_{\Psi_{N'}} = cN' \varepsilon^{-n/(s-1)} \log(\varepsilon^{-1})$ active weights,
    \end{enumerate}
    for the Decoder, and
    \begin{enumerate}
        \item[-] $L_{\Phi_n} = c \log(\varepsilon^{-1})$ layers,
        \item[-] $\omega_{\Phi_n} = cn \varepsilon^{-(p+q+1)} \log(\varepsilon^{-1})$ active weights,
    \end{enumerate}
    for the Deep Feed-Forward Neural Network.
\end{theorem}
\begin{proof}
    Similar as to the proof of Theorem (\ref{theo: PAC bound for POD-DL-ROM}), we will try to bound each quantity
    of the statement in Theorem (\ref{theo: relative error DOD+DNN}), but with the catch, that $\Ecal_{\text{DOD}}$ is
    stochastic in nature, and thus we will employ Remark (\ref{remark: discrete optimal DOD}) to find an auxilary term,
    which is independent of the data, to which the approximation of the neural net of the DOD is sufficiently "close". 
    
    The sampling error can be bounded in the same way as before, i.e. we can directly bound 
    $\Ecal_S = \Ecal_S(N_{s_1}, N_{s_2}, N_t)$ independently of $N'$, under the Assumption (\ref{assumption sample})
    by the Weak Law of Large Numbers, i.e. for all $1 > \delta > 0$ and $\varepsilon > 0$ there are $N_{s_1}, N_{s_2}, N_t$, such that
    \begin{equation} \label{eq: S error dod-dl-rom}
        \prob[\Ecal_S((N_{s_1}, N_{s_2}, N_t)) < \varepsilon / 4] > 1 - \delta.
    \end{equation}
    Now recall the definition of $\Ecal_{\text{DOD}}$ in the proof of Theorem (\ref{theo: relative error DOD+DNN}) and 
    further establish
    \begin{align*}
        \Ecal_{\text{DOD}} &= \frac{m^{-1} \vert \Theta \times \Theta' \times [0, T] \vert}{N_{data}^{1/2}} 
        \, \bigg\vert \, \sum_{i = 1}^{N_{s_1}} \sum_{j = 1}^{N_{s_2}} \sum_{k = 1}^{N_t} 
        \norm{u_h^{\mu_i, \nu_j, k \Delta t} - \V_{\mu_i, k \Delta t} \V_{\mu_i, k \Delta t}^T u_h^{\mu_i, \nu_j, k \Delta t}}^2 \, \bigg\vert^{1/2} \\
        &\leq m^{-1} \sup\limits_{(\mu, t) \in \Theta \times [0, T]} \, \bigg\vert \, 
        \frac{\vert \Theta \times \Theta' \times [0, T] \vert}{N_{s_2}}  \sum_{j = 1}^{N_{s_2}} 
        \norm{u_h^{\mu, \nu_j, t} - \V_{\mu, t} \V_{\mu, t}^T u_h^{\mu, \nu_j, t}}^2 \, \bigg\vert^{1/2}, 
    \end{align*}
    enabling us to use Proposition (\ref{prop: eigenvalues and optimal projection under sampling}) in order 
    to switch to the eigenvalues $\sigma(\mu, t)_k$ for $(\mu, t) \in \Theta \times [0, T]$ given by
    Equation (\ref{eq: definition virtual eigenvalues}) 
    and employing the Remark (\ref{remark: discrete optimal DOD}) with 
    the bound
    \begin{equation*}
        \frac{\varepsilon \, m^2}{4 \vert \Theta \times \Theta' \times [0, T] \vert^{1/2}},
    \end{equation*}
    which directly gives us the approximation result via the defition of $N'$ in the prelminaries of this section
    \begin{equation} \label{eq: dod error dod-dl-rom}
        \Ecal_{\text{DOD}} < \frac{\varepsilon}{4} + 
        m^{-1} \sup\limits_{(\mu, t) \in (\Theta \times [0, T])} 
        \sqrt{\sum_{k > N'} \sigma(\mu, t)_k^2} 
        \leq \frac{\varepsilon}{2}.
    \end{equation}

    \notdone{Include the cost via Yarotski}
    \notdone{Argue the actual part for the NN error}

    Moreover we can use the exact same arguments as in the bound for $\Ecal_{\text{NN}}$ for the POD-DL-ROM as in the
    proof of Theorem (\ref{theo: PAC bound for POD-DL-ROM}), which yields
    \begin{equation} \label{eq: nn error dod-dl-rom}
        \Ecal_{\text{NN}} < \frac{\varepsilon}{4}.
    \end{equation}
    Finally putting each Equation, i.e. (\ref{eq: S error dod-dl-rom}), (\ref{eq: dod error dod-dl-rom}) and (\ref{eq: nn error dod-dl-rom}),
    together, we get 
    \begin{equation*}
        \Ecal_R \leq \Ecal_S + \Ecal_{\text{DOD}} + \Ecal_{\text{NN}} < \frac{\varepsilon}{4} +  \frac{\varepsilon}{2}
        + \frac{\varepsilon}{4} = \varepsilon,
    \end{equation*}
    with probability greater than $1 - \delta$ concluding our proof.
\end{proof}






\section{Comparison between existing Techniques}